\subsection{匹配圈}\label{sec:16g}
令 $\pi:E\to S$ 为 Lefschetz 纤维化，仍照例假设正合。取嵌入路径 $\mu:[-1,1]\to\operatorname{int}(S)$，满足 $\mu^{-1}(\operatorname{Critv}(\pi))=\{-1,1\}$。可将它分为具有相同起点的一对消失路径 $\gamma_\pm(t)=\mu(\pm t)$，$t\in[0,1]$，因而得到一对消失圈 $V_{\gamma_\pm}\subset M=E_{\mu(0)}$。先看一个特别简单的情形，即两者相等。此时，
\begin{equation}\label{eq:16.8}
\Sigma_\mu=\Delta_{\gamma_-}\cup\Delta_{\gamma_+}
\end{equation}
是总空间 $E$ 的光滑 Lagrangian 子流形。由 Lefschetz 顶针的定义，对所有 $-1<t<1$，沿 $\mu$ 的平行移动把交 $\Sigma_\mu\cap E_{\mu(t)}$ 互相映到彼此，这在重叠处 $\Delta_{\gamma_-}\cap\Delta_{\gamma_+}=V_{\gamma_\pm}$ 附近给出局部坐标 $(-1,1)\times V_{\gamma_\pm}$。由于它是两个球沿边界粘合的结果，$\Sigma_\mu$ 必为同伦球面。若假设 $V_{\gamma_\pm}$ 的标架同痕，则它在微分意义下是标准球面。事实上，给定两个标架之间的同痕，可把\eref{16.1}的微分同胚合在一起，得到 $\Sigma_\mu$ 的标架。将\eref{16.8}称为\emph{朴素匹配圈构造}。

要求两个消失圈严格相等似乎很不现实，但可以稍微放宽假设，代价是失去对所得球面在 $E$ 中位置的控制。具体地，假设 $V_{\gamma_\pm}$ 在 $M$ 中作为带标架正合 Lagrangian 球面同痕。此时称 $\mu$ 为\emph{匹配路径}。两个消失圈之间的任意带标架正合 Lagrangian 同痕 $I=(i,j_0,j_1)$，也相应称为\emph{匹配同痕}。\znote{230}{匹配圈与消失圈的区别在于“控制权”：消失圈由纤维化完全决定，而匹配圈要先把两个消失圈在中心纤维里同痕到一起，再沿 $\mu$ 把两块顶针接起来，因此只能确定到 $E$ 中的同痕类，位置不受控制——这正是本节开头所说“失去对所得球面在 $E$ 中位置的控制”。$I=(i,j_0,j_1)$ 中的 $i^s$ 是把 $V_{\gamma_-}$ 同痕到 $V_{\gamma_+}$ 的一族 Lagrangian 嵌入，$j_0,j_1$ 则把两端的标架对齐；所以 $I$ 既是位置的粘合数据，也是标架的粘合数据，舍去它 $\Sigma_\mu$ 就不再是带标架球面。[来源：原书 §(16g)，印刷页 230。]}初等辛几何表明，$\operatorname{Aut}(M,\partial M)$ 中存在同痕 $(\phi^s)$，使 $\phi^0=\Id$、$\phi^s\circ i^0=i^s$。由\stref{引理}{15.3}，存在给定正合 Lefschetz 纤维化的变形，沿纤维恒定并集中在 $E_{\mu(-1/4)}$ 的某个邻域，使沿 $\mu|[-1/2,0]$ 的平行移动映射相应变形为
\[
h_{\mu|[-1/2,0]}^s=\phi^s\circ h_{\mu|[-1/2,0]};
\]
右边的平行移动是初始结构，即 $s=0$ 时的平行移动。对消失圈的影响为 $V_{\gamma_-}^s=\phi^s(V_{\gamma_-})$，而 $V_{\gamma_+}$ 保持不变。在变形的另一端 $s=1$，得到一个两个消失圈相等的正合 Lefschetz 纤维化结构，因而可用朴素匹配圈构造在总空间中定义
\SourcePage{239}
Lagrangian 球面。由于在 $i^0,i^1$ 对有关向量空间给出的认同下，消失圈的两个标架也相同，朴素匹配圈便带有标架。对反向变形应用\stref{引理}{7.1}，得到从 $(E,\omega_E^1)$ 到 $(E,\omega_E=\omega_E^0)$ 的共形辛嵌入 $\rho$，并用它把朴素匹配圈映回原来的辛形式。在带标架正合 Lagrangian 球面的同痕意义下，结果不依赖于 $(\phi^s)$、Lefschetz 纤维化变形以及该嵌入 $\rho$ 的选择。称其为与 $\mu,I$ 关联的\emph{匹配圈}，记为 $\Sigma_{\mu,I}\subset E$。为简化记号，在有明显的同痕选择，或不关心采用哪一种选择时，仍写作 $\Sigma_\mu$。

\StatementNumber{11}\begin{note}\label{st:16.11}
同痕 $I$ 在构造中的作用不易估量。保持端点不动地变形 $I$，不影响匹配圈的同痕类。在最低维情形 $n=1$，曲面 $M$ 中正合嵌入圆周所成的空间具有平凡的 $\pi_1$，故 $\Sigma_{\mu,I}$ 与 $I$ 无关。更高维中的相应问题似乎很难回答。不过，后面结合\stref{命题}{18.21}与\stref{命题}{18.13}将看到，$\Sigma_{\mu,I}$ 作为 Donaldson–Fukaya 范畴对象的同构类与 $I$ 无关。
\end{note}

一般而言，无法保证 $E$ 中给定的 Lagrangian 球面能表示为匹配圈；不过，Lefschetz 束背景下此类渐近结论可参看\bcite{16}。此外，不同同痕类的匹配路径可以给出同痕的匹配圈。这不是偶发现象，而是任何足够复杂的 Lefschetz 纤维化中都可以预期出现的情况。下例出自 Auroux：

\StatementNumber{12}\begin{example}\label{st:16.12}
为简单起见，考虑最低维情形 $\dim(E)=4$。假设底空间 $S$ 中有两条嵌入路径 $\mu,\nu$，各自连接一对临界值，四个临界值两两不同；两条路径恰相交一次，显然是在各自的某个内点，并且横截相交。假设 $\mu$ 为匹配路径，而 $\nu$ 两半对应的消失圈横截相交且只有一个交点。令 $t_\nu$ 为沿 $\nu$ 的半扭转。由于沿相交一次的圆周的 Dehn 扭转满足辫关系，$\mu'=t_\nu^3(\mu)$ 仍为匹配路径，且相应的 Lagrangian 二维球面 $\Sigma_{\mu'}$ 与 $\Sigma_\mu$ 同痕。一种直接但稍不直观的做法是：取所涉及全部 Lefschetz 顶针的邻域；从胞腔附着的观点考察它；认出它与 $T^*S^2$ 的某个紧部分辛同构；再应用\bcite{74}，证明 Lagrangian 二维球面在同痕意义下的唯一性。这一构造还有一个略微不那么有趣的变体：假设 $\nu$ 两半对应的消失圈不交，并取 $\mu'=t_\nu^2(\mu)$。
\end{example}
